\documentclass[11pt,a4paper]{article}
\usepackage[utf8]{inputenc}
\usepackage[T1]{fontenc}
\usepackage{lmodern}
\usepackage[margin=1in]{geometry}
\usepackage{amsmath,amssymb,amsthm}
\usepackage{booktabs}
\usepackage{array}
\usepackage{microtype}
\usepackage[hidelinks]{hyperref}

\newtheorem{theorem}{Theorem}
\newtheorem{corollary}[theorem]{Corollary}
\newtheorem{proposition}[theorem]{Proposition}
\theoremstyle{remark}
\newtheorem{remark}[theorem]{Remark}

\newcommand{\tier}[1]{{\small\textsf{[#1]}}}
\newcommand{\src}[1]{{\footnotesize\texttt{#1}}}
\newcommand{\LamDH}{\Lambda_{\mathrm{DH}}}
\newcommand{\PhiF}{\Phi_F}
\newcommand{\HDH}{H^{\mathrm{DH}}}

\title{\bfseries A Two-Sided Certified Bracket for the de Bruijn--Newman Constant
of the Davenport--Heilbronn Function, and What the Heat Flow Does Not Transfer}

\author{Elias De Jes\'us\\[2pt]
{\small Independent Researcher}\\[1pt]
{\small ORCID \href{https://orcid.org/0009-0007-0190-9143}{0009-0007-0190-9143}}\\[1pt]
{\small \texttt{dejesuselias10@gmail.com}}}

\date{October 2026 --- Technical Note\\[2pt]
{\small\texttt{doi:10.5281/zenodo.23628224}}}

\begin{document}
\maketitle

\begin{abstract}
\noindent
Let $\PhiF(x)=4e^{3x}\sum_{n\ge1}na_ne^{-\pi n^2e^{4x}/5}$ be the kernel of the completed
Davenport--Heilbronn function and $\HDH_\tau(z)=\int_0^\infty e^{\tau x^2}\PhiF(x)\cos(zx)\,dx$
its heat-flow family, so that $\HDH_0(z)=\tfrac12\Lambda_F(\tfrac12+iz/2)$. Write $\LamDH$
for the de Bruijn--Newman constant of this family: all zeros of $\HDH_\tau$ are real exactly
when $\tau\ge\LamDH$. We prove
\[
  6.8268151773112564447\times10^{-29}\ \le\ \LamDH\ \le\ 1.602537483559822777787819 ,
\]
the upper end from de Bruijn's strip theorem applied to the elementary zero-free abscissa of
$\Lambda_F$, the lower end from a Rouch\'e persistence argument certified in ball arithmetic
on one disc about the first off-line zero. Both ends are numerical, and the bracket is
nonempty: the Davenport--Heilbronn flow has a genuinely positive transition parameter. The
second half of the note states what this does \emph{not} give. De Bruijn's theorem needs
only conjugate reflection, rapid decay and a bounded initial strip; the Davenport--Heilbronn
function has all three, together with a \emph{positive} kernel, while its finite-window Weil
form at $\tau=0$ is certifiably indefinite. So no argument resting on that package alone can
carry positivity backward to $\tau=0$ for $\zeta$, and we record the three properties a
hypothesis must have to escape the obstruction.
\end{abstract}

\section*{0.\quad Non-claims}
\addcontentsline{toc}{section}{Non-claims}

Stated first.

\begin{enumerate}\itemsep3pt
\item \textbf{Nothing about RH or GRH.} Every statement below is about a deformation family
and about what a finite-window detector built from its zeros can be made to do. None is
about where the zeros of $\zeta$ are.

\item \textbf{The known bounds $0\le\Lambda\le0.22$ for $\zeta$ are restated, not
improved}, and nothing here bears on them. $\Lambda\ge0$ is Rodgers--Tao~\cite{RT};
$\Lambda\le0.22$ is Polymath15~\cite{P15}; $\Lambda\le\tfrac12$ is de Bruijn's classical
value~\cite{DB50}. The bracket of this note is for the Davenport--Heilbronn family only, and
its $\zeta$ analogue is exactly those two results and nothing from here.

\item \textbf{Nothing about whether heat flow can prove RH}, beyond the scope statement of
Section~\ref{sec:transfer}, which is a statement about one package of hypotheses and not
about a programme.

\item \textbf{The quasi-Riemann-hypothesis preprints of October 2026 are not cited and not
used.} No statement here depends on them, and none is asserted about them.

\item \textbf{Tiers.} The upper bound is \tier{T1} \emph{conditional on} the strip bound it
consumes, which is itself \tier{T1} in the companion extension. The lower bound is
\tier{T2-C}: $m$, $C$ and the winding number are ball enclosures. Every zero
\emph{location} is \tier{T2} --- a numerical input, not independently certified here --- and
the disc centre is one of those, though Section~\ref{sec:lower} shows nothing depends on it.

\item \textbf{The bracket is not a localisation.} The upper end is a \emph{sufficient}
condition from the strip theorem, not a computed transition point, and the lower end does
not say that $\LamDH$ is small. Where $\LamDH$ actually sits in this interval is not
determined here and no attempt is made on it.

\item \textbf{$\varepsilon\approx 7\times10^{-29}$ is quantitatively weak and is not
presented otherwise.} Section~\ref{sec:lower} explains why the order of magnitude is
structural rather than a tuning artefact, and the disc and the bound were not loosened to
produce a larger number.

\item \textbf{No novelty claim.} De Bruijn's theorem, the heat equation and the Rouch\'e
argument are classical. The only assembly the author has not seen elsewhere is the
application of the strip theorem to the Davenport--Heilbronn kernel, and no priority is
claimed for it.
\end{enumerate}

\section*{Epistemic tiers}
\addcontentsline{toc}{section}{Epistemic tiers}

\begin{center}
\begin{tabular}{@{}lll@{}}
\toprule
tier & meaning & here \\
\midrule
\tier{T1} & proved, by an argument written out & the kernel identity, the strip transfer \\
\tier{T2-C} & interval/ball certified & $m$, $C$, the winding number, $\varepsilon$ \\
\tier{T2-N} & numerical, high precision, not certified & the identity cross-checks \\
\tier{T2} & a numerical input quoted as such & zero locations, the disc centre \\
\tier{T3} & reasoned judgement & the scope statements of \S\ref{sec:transfer}--\S\ref{sec:filter} \\
\bottomrule
\end{tabular}
\end{center}

\section{Setup}\label{sec:setup}

\paragraph{The family.} Let $F(s)=\tfrac12[(1-i\kappa)L(s,\chi)+(1+i\kappa)L(s,\bar\chi)]$ be
the Davenport--Heilbronn function~\cite{DH36,BG11}, $\chi$ the odd primitive character
modulo $5$ and $\kappa$ the constant making the functional equation exact, and let
$\Lambda_F$ be its completion. From the Mellin representation
$\Lambda_F(s)=\int_0^\infty\omega(t)t^{(s+1)/2}\,dt/t$ with
$\omega(t)=\sum_{n\ge1}na_ne^{-\pi n^2t/5}$, the substitutions $t=e^{4x}$,
$s=\tfrac12+iz/2$ give $\Lambda_F(\tfrac12+iz/2)=\int_{\mathbb R}\PhiF(x)e^{izx}dx$ with
\begin{equation}\label{eq:kernel}
  \PhiF(x)=4e^{3x}\sum_{n\ge1}na_n\exp(-\pi n^2e^{4x}/5),\qquad
  \HDH_\tau(z)=\int_0^\infty e^{\tau x^2}\PhiF(x)\cos(zx)\,dx,
\end{equation}
so that $\HDH_0(z)=\tfrac12\Lambda_F(\tfrac12+iz/2)$. The derivation is \tier{T1} and is
written out in the audit record~\cite[\S2]{audit}; it is not inferred from the $\zeta$ case.
Coordinates are kept separate throughout: $s$ is the Dirichlet variable, $z=-2i(s-\tfrac12)$
the spectral coordinate --- so for $s=\beta+i\gamma$, $\operatorname{Im}z=-2(\beta-\tfrac12)$
and \emph{the critical line is the real axis} --- and $\tau$ is the deformation parameter,
which is not time and carries no mechanism.

\paragraph{The hypotheses of Theorem~\ref{thm:dB}, checked one by one.}
Each is established, not inherited.
\begin{itemize}\itemsep2pt
\item \emph{$\PhiF$ real.} \tier{T1} The Davenport--Heilbronn coefficients are
$a_n=\operatorname{Re}\chi(n)+\kappa\operatorname{Im}\chi(n)$ with $\kappa$ real, hence real;
so \eqref{eq:kernel} is a real series.
\item \emph{$\PhiF$ even.} \tier{T1} $\Lambda_F$ satisfies
$\Lambda_F(s)=\Lambda_F(1-s)$~\cite[\S1]{dhrep}, which in $s=\tfrac12+iz/2$ is $z\mapsto-z$.
So $\int\PhiF(x)e^{izx}dx$ is even in $z$, hence equals $\int\PhiF(-x)e^{izx}dx$, and by
uniqueness of Fourier transforms $\PhiF(-x)=\PhiF(x)$. Evenness is therefore
\emph{equivalent} to the functional equation, and the numerical check below is a check of
the substitution, not the source of the fact.
\item \emph{Decay.} \tier{T1} $\PhiF(x)=O\bigl(e^{3|x|}e^{-(\pi/5)e^{4|x|}}\bigr)$ is doubly
exponential, so $|\PhiF(t)|=O(\exp(-|t|^{\gamma}))$ for \emph{every} $\gamma$, and in
particular for some $\gamma>2$.
\item \emph{The strip.} \tier{T1} given $\sigma_0$; this is \S\ref{sec:upper}.
\end{itemize}

\noindent \tier{T2-N} As cross-checks, the kernel \eqref{eq:kernel} reproduces
$\Lambda_F(\tfrac12+iz/2)$ to worst relative residual $5.10986\times10^{-20}$ and is even to
$3.43924\times10^{-29}$. The flow equation
$\partial_\tau H=-\partial_z^2H$ was checked by finite differences at
$2.76562\times10^{-22}$ and $5.17887\times10^{-22}$, and the same checks on the $\zeta$
kernel reproduce its normalisation to $7.17465\times10^{-43}$. The flow equation plays no
part in Theorem~\ref{thm:dB} or in either bound; it is recorded because it identifies the
family as the heat flow of the sources.

\paragraph{Positivity, recorded but not needed.} With $T=\pi e^{4x}/5\ge\pi/5$,
$|a_n|\le1$ and $a_1=1$, factoring out the leading term gives
\begin{equation}\label{eq:pos}
  \sum_{n\ge1}na_ne^{-n^2T}\ \ge\ e^{-T}\Bigl[1-\sum_{n\ge2}ne^{-(n^2-1)T}\Bigr],
\end{equation}
and the bracket is \emph{increasing} in $T$, so one evaluation bounds it below for all
$x\ge0$: it is $0.6763199240515301$ at $T=\pi/5$. With evenness, $\PhiF>0$ on
$\mathbb R$ (sampled minimum $5.6615196\times10^{-44408}$ on $[0,3]$). The unfactored
difference must not be used in its place, since it tends to $0$ as $T$ grows; an earlier
version of the record did use it and the correction is logged in Appendix~\ref{app:retr}.
Positivity is reported because the surrounding framework of~\cite{NW} assumes a positive
measure, so the Davenport--Heilbronn case sits inside that framework and not merely inside
the single theorem. \emph{The theorem itself does not use it}, which is the point of
\S\ref{sec:transfer}.

\paragraph{The constant, and the theorem.} The family has a well-defined de Bruijn--Newman
constant $\LamDH$: all zeros of $\HDH_\tau$ are real exactly when $\tau\ge\LamDH$, by the
same three facts that give $\Lambda$ for $\zeta$ --- forward preservation, an upper member
of the set, and boundedness below~\cite{NW,Ne76}. We use de Bruijn's theorem in the
normalisation reconciled in~\cite[\S S]{audit}:

\begin{theorem}[de Bruijn {\cite[Thm.~13]{DB50}}, in the form stated and proved as
{\cite[Thm.~1]{Ne76}}]\label{thm:dB}
\tier{T1} Let $f$ be a real, even, integrable function on $\mathbb R$ with
$|f(t)|=O(\exp(-|t|^{\gamma}))$ as $|t|\to\infty$ for some $\gamma>2$, and suppose all the
zeros of the Fourier transform of $f$ lie in the strip $|\operatorname{Im}z|\le\Delta$. Then
for $\delta\ge0$ all the zeros of the Fourier transform of $\exp(\delta t^2)f(t)$ lie in
\[
  |\operatorname{Im}z|\ \le\ \bigl[\max(\Delta^2-2\delta,\,0)\bigr]^{1/2}.
\]
\end{theorem}

\noindent \textbf{This is the statement in our own multiplier convention}, $\exp(\delta
t^2)$, so all zeros are real once $\delta\ge\Delta^2/2$ with no further bookkeeping. We quote
Newman rather than the survey restatement~\cite[Thm.~7]{NW}, whose multiplier is
$e^{\lambda t^2/2}$ and which therefore carries a factor $\lambda=2\delta$ relative to ours.
\textbf{Positivity of $f$ is not a hypothesis of this theorem} and is not used for the upper
bound; it is a hypothesis of the surrounding framework of~\cite{NW}, and the bracket
\eqref{eq:pos} above belongs to the Bochner argument of \S\ref{sec:filter}, not here.

\paragraph{The $z$-convention, stated because it is a known hazard.}
\tier{T3} The threshold $\Delta^2/2$ is a statement about the half-width of the strip
\emph{in the variable the transform is taken in}, and the literature uses at least three
such variables for $\zeta$, so the same theorem appears in print with three different
numbers. We therefore fix ours and check it rather than inherit a value.

Here $z=-2i(s-\tfrac12)$, so $\operatorname{Im}z=-2(\beta-\tfrac12)$ and the $z$-strip
half-width is \emph{twice} the $s$-strip half-width $\Delta_s$. With
$\Delta=2\Delta_s$ the threshold is $2\Delta_s^2$. For $\zeta$, $\Delta_s=\tfrac12$ gives
$\Delta=1$ and $\tau\ge0.5$ --- the value quoted by Polymath15 and by Rodgers--Tao, and the
reason that reproduction is worth printing: it pins the convention. For comparison,
Newman~\cite[p.~246]{Ne76} works in $z=-i(s-\tfrac12)$, where $\zeta$'s strip has
$\Delta=\tfrac{1}{2}$ and the same theorem gives $\tfrac{1}{8}$; de Bruijn's own normalisation
is a third, $\Xi(2z)=\int\varphi(t)e^{izt}dt$. None of these disagree: they are one theorem
in three variables. In the Titchmarsh-type scaling the bound of
Proposition~\ref{prop:upper} would read $\Delta_s^2/2=0.40063437088995569$, exactly a
quarter of our value; and Polymath15's $0.22$ is in \emph{our} convention, so it is
comparable with $0.5$ and not with $\tfrac{1}{8}$.

\section{The upper bound}\label{sec:upper}

\begin{proposition}\label{prop:upper}
\tier{T1} \emph{given} the strip bound.
$\LamDH\le2(\sigma_0-\tfrac12)^2=1.602537483559822777787819$, where
$\sigma_0=1.395136158235109721061359$.
\end{proposition}

\begin{proof}
Every zero of $\Lambda_F$ satisfies $|\beta-\tfrac12|\le\sigma_0-\tfrac12$, i.e.
$|\beta-\tfrac12|\le0.895136158235109721061359$ \cite[\S3]{dhnote}. Since
$\operatorname{Im}z=-2(\beta-\tfrac12)$, the zeros of $\HDH_0$ lie in
$|\operatorname{Im}z|\le\Delta_{\mathrm{DH}}=2(\sigma_0-\tfrac12)
=1.790272316470219442122718$. Theorem~\ref{thm:dB} with $\lambda=2\tau$ then puts every
zero on the real axis once $\tau\ge\Delta_{\mathrm{DH}}^2/2$, which is the stated number.
\end{proof}

Three things about this bound should be said plainly, because each is a limitation.

\emph{It is only as good as the strip.} The input $\sigma_0$ is an \emph{elementary
coefficient estimate} --- the proved zero-free abscissa of $\Lambda_F$, obtained by bounding
the Dirichlet series against its first term --- not a deep fact about the
Davenport--Heilbronn function. A sharper strip lowers the bound quadratically, and a
hypothetical $\sigma_0\to\tfrac12$ would drive it to $0$. Nothing here attempts to sharpen it.

\emph{It is sufficient, not sharp.} Proposition~\ref{prop:upper} names a $\tau$ by which all
zeros are certainly real. It does not locate the transition, and the true $\LamDH$ may be far
below.

\emph{The only arithmetic it consumes is that one number.} This is the observation that
Sections~\ref{sec:transfer}--\ref{sec:filter} turn into an obstruction: the toy family
$z^2+a^2-2\tau$, which has no arithmetic content whatever, consumes exactly as much and
saturates Theorem~\ref{thm:dB} exactly~\cite[\S5]{audit}.

\section{The lower bound, certified}\label{sec:lower}

The lower end is a local argument on a single disc. It uses no flow numerics, tracks no
zero, and does not differentiate the detector family.

\begin{theorem}\label{thm:lower}
\tier{T2-C} $\LamDH\ \ge\ \varepsilon\ \ge\ 6.8268151773112564447\times10^{-29}$.
\end{theorem}

\begin{proof}
Let $D$ be the closed disc of radius $r$ about $z_1$, the first off-line zero of $\HDH_0$ in
the spectral coordinate, and suppose $\overline D$ misses the real axis. Put
\[
  m=\min_{\partial D}|\HDH_0|,\qquad M=\max_{\partial D}|\operatorname{Im}z|,\qquad
  C=\int_0^\infty x^2|\PhiF(x)|e^{\tau_0x^2+Mx}\,dx .
\]
For $|\tau|\le\tau_0$ and $z\in\partial D$,
\begin{equation}\label{eq:pert}
  |\HDH_\tau(z)-\HDH_0(z)|
  =\Bigl|\int_0^\infty(e^{\tau x^2}-1)\PhiF(x)\cos(zx)\,dx\Bigr|\ \le\ |\tau|\,C ,
\end{equation}
by the single inequality
$|e^{\tau x^2}-1|=\bigl|\int_0^{\tau x^2}e^t\,dt\bigr|\le|\tau|x^2e^{|\tau|x^2}
\le|\tau|x^2e^{\tau_0x^2}$, valid for $|\tau|\le\tau_0$, together with
$|\cos(zx)|\le\cosh(x\operatorname{Im}z)\le e^{Mx}$ for $x\ge0$. If $|\tau|\,C<m$ then
$|\HDH_\tau-\HDH_0|<|\HDH_0|$ on $\partial D$, so by Rouch\'e $\HDH_\tau$ has the same
number of zeros in $D$ as $\HDH_0$. That number is nonzero (established below), so
$\HDH_\tau$ has a zero in $D$, which is \emph{nonreal} because $\overline D$ misses the real
axis; hence not all zeros of $\HDH_\tau$ are real, hence $\tau<\LamDH$. This holds for every
$|\tau|\le\varepsilon:=\min(\tau_0,m/(2C))$, so $\LamDH\ge\varepsilon$. The four certified
inputs are tabulated in Appendix~\ref{app:cert}; substituting them gives the stated value.
\end{proof}

\paragraph{The disc, and that its closure misses the real axis.} \tier{T2-C} In the spectral
coordinate the critical line is the real axis and the first off-line zero sits at
$z_1=2\gamma_1-2i\delta_1$ with
$\operatorname{Re}z_1=171.398696970755184343858535418$ and
$\operatorname{Im}z_1=-0.617034364913274771106703921214$. With $r=0.3$,
\[
  |\operatorname{Im}z_1|-r=0.31703436491327477111>0,
\]
so $\overline D$ is disjoint from the real axis, with that margin.

\paragraph{$m>0$, by arc enclosures.} \tier{T2-C} $\partial D$ is partitioned into
$3000$ arcs of angular width $2\pi/3000$. Each arc
$\{c+re^{i\theta}:|\theta-\theta_k|\le\Delta\theta/2\}$ lies in the ball of centre
$c+re^{i\theta_k}$ and radius $2r\sin(\Delta\theta/4)=0.000314159251004$, because
$|re^{i\theta}-re^{i\theta_k}|=2r|\sin((\theta-\theta_k)/2)|$; and
$\HDH_0=\Lambda_F(\tfrac12+iz/2)/2$ is enclosed on each such ball using ball arithmetic for
$\Gamma$ and the Hurwitz $\zeta$~\cite{Arb}. The worst relative enclosure radius over the
arcs is $0.064946579$, and
\[
  m\ \ge\ 6.5935981835929959399\times10^{-30}\ >\ 0 ,
\]
a plain decimal floor: the enclosure's lower endpoint rounded \emph{down}, so the inequality
is exact. The partition size was \emph{measured, not guessed}. The Hurwitz-$\zeta$ sum
amplifies an input ball radius by a factor of about $106$ at this height --- that is the
genuine size of the derivative, not an artefact of the arithmetic --- so a $60$-arc
partition gives a relative enclosure radius of $3.30$ and fails to decide anything, while
$3000$ arcs give $0.0649$.

\paragraph{The zero inside $D$ is certified, not assumed.} \tier{T2-C} Summing the principal
argument increments of the arc enclosures,
\[
  \frac{1}{2\pi}\sum\Delta\arg = 1\ \ (\pm\,3\times10^{-20}),\qquad
  \max_k\bigl(|\Delta_k\arg|+\text{both arc uncertainties}\bigr)=0.1313552700\ <\ \pi .
\]
Every increment is therefore unambiguous and the winding number is \emph{exactly} $1$: the
function $\HDH_0$ has exactly one zero in $D$, counted with multiplicity. \textbf{So the
census is not an input to Theorem~\ref{thm:lower}.} The disc centre came from it, but
nothing in the proof requires that centre to be a zero, and the \tier{T2} status of the
zero locations does not propagate into the bound.

\paragraph{$M$ and $C$.} \tier{T2-C} $M=|\operatorname{Im}z_1|+r=0.91703436491327477111$ and
$\tau_0=0.1$. The integral $C$ is split at $x=2.000$. On $[0,x_{\mathrm{cut}}]$ it is
evaluated by rigorous ball integration of the explicit integrand, with the kernel truncated
at $n\le40$ and the omitted terms carried as an error ball, giving
$0.048291904880526492445$. For the tail, $|\PhiF(x)|\le8e^{3x}e^{-\pi e^{4x}/5}$ for
$x\ge1$, so with $f(x)=x^2\cdot8e^{3x}e^{-\pi e^{4x}/5}e^{\tau_0x^2+Mx}$,
\[
  (\log f)'(x)=\frac2x+3-\frac{4\pi}{5}e^{4x}+2\tau_0x+M\ \le\ -1
  \qquad (x\ge x_{\mathrm{cut}}),
\]
since $(4\pi/5)e^{8}=7491.9646$ dominates the remaining terms there and grows far faster.
Hence $f(x)\le f(x_{\mathrm{cut}})e^{-(x-x_{\mathrm{cut}})}$ and
$\int_{x_{\mathrm{cut}}}^\infty f\le f(x_{\mathrm{cut}})=4.481699445\times10^{-809}$. So
$C\le0.048291904880526492446$, rounded \emph{up}.

\begin{remark}[why $\varepsilon$ is this small, and what that does and does not mean]
\label{rem:small}
\tier{T3} The order of magnitude is structural. $m\approx 6.6\times10^{-30}$ is small because
$\Lambda_F$ carries a gamma factor whose modulus at this height is of that order; that is a
\emph{prefactor}, present for any disc about this zero, not a cancellation. Meanwhile
$C\approx 0.048$ is $O(1)$ because the perturbation bound \eqref{eq:pert} replaces
$\cos(zx)$ by $e^{Mx}$ and so discards every bit of the oscillatory cancellation that makes
$\HDH_0$ small in the first place. The ratio inherits the gamma decay on one side and none
of it on the other, so $\varepsilon\sim 10^{-29}$ for any reasonable radius and no choice of
$r$ changes the exponent. Improving it would need a perturbation estimate that keeps the
oscillation --- a saddle-point treatment of
$\int(e^{\tau x^2}-1)\PhiF(x)\cos(zx)\,dx$ --- which is a different and more delicate
undertaking, and was not attempted. What the bound settles is the qualitative question with
a certificate rather than a continuity appeal: $\LamDH$ is bounded away from zero.
\end{remark}

\section{The bracket, and the $\zeta$ comparison}\label{sec:bracket}

Combining Propositions~\ref{prop:upper} and Theorem~\ref{thm:lower}:
\begin{equation}\label{eq:bracket}
  \boxed{\ 6.8268151773112564447\times10^{-29}\ \le\ \LamDH\ \le\
  1.602537483559822777787819\ }
\end{equation}
with the lower end \tier{T2-C} and the upper end \tier{T1} given the strip bound.

The same two-sided reasoning applied to $\zeta$ returns the known interval
$0\le\Lambda\le0.22$: the lower end is a theorem --- Rodgers--Tao~\cite{RT}, and the
quantitative lower bounds that preceded it were obtained by a quite different route, the
Lehmer-pair repulsion argument of Csordas--Smith--Varga~\cite{CSV94} rather than by local
persistence --- and the upper end is the Polymath15 computation~\cite{P15}, sharpening de
Bruijn's $\tfrac12$. The structural difference is the one worth naming. For $\zeta$ the
margin has collapsed onto the single point $\Lambda=0$, which \emph{is} the Riemann
hypothesis, so the interval cannot be narrowed from the left at all without proving it. For
the Davenport--Heilbronn function the margin is open, and for a reason available at the
start: the function has nonreal zeros at $\tau=0$, so persistence gives a strictly positive
lower bound with no arithmetic input beyond the location of one of them. The bracket
\eqref{eq:bracket} is therefore easy in a way the $\zeta$ bracket is not, and that is the
content of the comparison --- not any suggestion that the method transfers.

\section{What the heat flow transfers, and what it does not}\label{sec:transfer}

\paragraph{The transfer, which is real.} \tier{T1} Let $\Psi_{ij}(z)=
\tfrac12[G_i(iz/2)G_j(-iz/2)+G_j(iz/2)G_i(-iz/2)]$, which at $z=2\gamma$ is exactly the
symmetrised explicit-formula kernel of the finite-window Weil form~\cite{deltaL3}; the
coordinate dictionary was checked against the existing implementation to
$1.11211\times10^{-40}$. For real $g$ and \emph{real} $z$ one has
$G_j(-iz/2)=\overline{G_j(iz/2)}$, so with $v_j=G_j(iz/2)$,
\[
  \Psi(z)=\operatorname{Re}(vv^*)=(\operatorname{Re}v)(\operatorname{Re}v)^T
  +(\operatorname{Im}v)(\operatorname{Im}v)^T\ \succeq\ 0 ,
\]
a sum of positive semidefinite rank-one terms. (In a pure parity sector one term vanishes
identically and the observed rank is $1$; the numerics caught a draft that said $2$, and the
conclusion is unaffected. Least eigenvalues at the precision floor,
$-2.296270406\times10^{-41}$, $-9.817490382\times10^{-42}$,
$-4.791060636\times10^{-41}$.) Hence if every zero of $H_\tau$ is real then
$Q_\tau=\sum_k m_k\Psi(z_k)$ is a sum of positive semidefinite matrices, and so
\[
  Q^\zeta_\tau\succeq0\ \ \text{for}\ \tau\ge0.5,\qquad
  Q^{\mathrm{DH}}_\tau\succeq0\ \ \text{for}\ \tau\ge1.602537483559822777787819,
\]
unconditionally, on \emph{every} finite window and every test space. That is a genuine
source-derived statement about the detector, and by Polymath15 the $\zeta$ threshold may be
replaced by $0.22$. The bridge does not fail at its first step.

\paragraph{And the obstruction.} \tier{T3} It points the wrong way. The hypotheses
Theorem~\ref{thm:dB} actually uses are conjugate reflection, rapid decay and a bounded
initial strip --- not kernel positivity, which belongs to the surrounding framework
and not to the theorem~\cite[\S S3]{audit}. The Davenport--Heilbronn function satisfies all
three \emph{and} has a positive kernel, with an explicit
$\Delta_{\mathrm{DH}}=1.790272316470219442122718$, while at $\tau=0$ its finite-window form
is certifiably indefinite: $\lambda_{\min}(Q^{\mathrm{DH}}_0)=
-0.679071346441480590986377586093\ldots$ on the $L=2$ odd level-$2$ space, enclosed in a
ball of radius $1.39\times10^{-116}$ \cite{dhnote}. Therefore:

\begin{quote}
\textbf{No argument resting on reflection symmetry, rapid decay, a bounded strip, kernel
positivity and heat evolution alone can transport positivity backward to $\tau=0$} ---
because every one of those hypotheses holds for the Davenport--Heilbronn function, whose
conclusion at $\tau=0$ is false.
\end{quote}

This is a statement about a package of hypotheses, and it is worth being exact about what it
leaves untouched. It does \emph{not} say that the arithmetic is destroyed by the
deformation: the coefficients $a_n$ are still present in
$\Phi_{F,\tau}=e^{\tau x^2}\PhiF$ at every $\tau$, and an earlier and sharper-sounding
slogan asserting otherwise is withdrawn (Appendix~\ref{app:retr}, point~3). What is missing
is an \emph{estimate} that uses the relations distinguishing $\zeta$ from
Davenport--Heilbronn. None is established here, and no claim is made about the
literature: the audit deliberately surveyed no recent work in this direction, so the
absence recorded here is an absence in this note and nowhere else. Nor has it been shown that the deformed family retains a usable Euler
product, gamma-factor decomposition or prime-side explicit formula for $\tau\ne0$; that is
an open question, not an impossibility.

\paragraph{Which links must fail, named.} \tier{T3} A backward comparison inequality
$\eta'(\theta)\le a(\theta)(1-\eta(\theta))$ propagating $\eta\le1$ from $\tau\ge\tau_{\mathrm{DH}}$
down to $\tau=0$ would contradict the certified witness. So somewhere on that path one of
its ingredients must fail, and two are identifiable. First, the channel representation
requires the standing assumption $B_\tau\succ0$ at every $\tau$, which is assumed and not
established along the path. Second, zero-tracking can meet collisions: differentiating the
complete sum term by term needs a uniform bound on $\sum_k|\Psi'(z_k)z_k'|$ with
$z_k'=H''/H'(z_k)$, which is not controlled near near-collisions, and a fixed contour can be
crossed as $\tau$ varies. The missing lemma is a uniform-in-$\tau$ tail bound together with
a contour zero-free throughout the interval; it is recorded as missing in the audit's
feasibility gate~\cite[\S3]{audit}, and for that reason no $\tau$-sweep was run and none is
reported. Neither of the two decisive statements of this note depends on it.

\section{Two geometries, and a filter}\label{sec:filter}

\paragraph{A positivity that is real and does not help.} \tier{T2-N} Because $\PhiF\ge0$, for
any nonnegative integrable $\Phi$ with transform $\widehat\Phi$ and any $c\in\mathbb C^n$,
\begin{equation}\label{eq:bochner}
  \sum_{j,k}c_j\overline{c_k}\,\widehat\Phi(y_j-y_k)
  =\int\Phi(x)\Bigl|\sum_jc_je^{-iy_jx}\Bigr|^2dx\ \ge\ 0 ,
\end{equation}
which is Bochner's theorem in one line. For the Davenport--Heilbronn kernel
$\widehat{\PhiF}(y)=\Lambda_F(\tfrac12+iy/2)$, so the matrix
$K_{jk}=\Lambda_F(\tfrac12+i(y_j-y_k)/2)$ \emph{must} be positive semidefinite. At
$y=(0,1,2.5,4,7)$ its five eigenvalues are
\[
  0.00002411328132,\quad 0.002983011492,\quad 0.1416783233,\quad
  1.363815621,\quad 5.657459774 ,
\]
positive semidefinite as required. And the same function's zero-based window form is
certifiably negative. There is no contradiction: these are different quadratic forms built
from different aspects of one function, one from the kernel's correlations and one from the
zeros. The nonimplication is worth stating plainly, because it is the trap a layered
programme walks into:

\begin{quote}
\textbf{Positivity in a source's Fourier-correlation geometry does not imply positivity in
its zero geometry.}
\end{quote}

So the missing bridge is not ``find something positive''. It is ``prove that \emph{this}
positive structure controls \emph{that} spectral form'', which is strictly harder.

\paragraph{The filter.} \tier{T3} Collecting the two preceding sections into a usable test:

\begin{quote}
\textbf{Statement F.} Any proposed argument for finite-window Weil positivity must contain
an independently established hypothesis that excludes the Davenport--Heilbronn function.
Reflection symmetry, rapid decay, a bounded zero strip, kernel positivity and heat evolution
do not supply one, because Davenport--Heilbronn has all of them and is certifiably
indefinite at $\tau=0$.
\end{quote}

For a candidate hypothesis to count, three things must hold, and the third is the one most
often skipped: it must be (i) \emph{independently proved} for the target function, not
assumed as part of the conclusion; (ii) \emph{false} for the control, so that it actually
separates them; and (iii) \emph{used in an estimate that controls zero geometry}, rather
than appearing only in the statement of the result. The one hypothesis the companion work
supplies that satisfies (i) and (ii) is the Euler product --- $\zeta$ has one and
Davenport--Heilbronn does not --- and (iii) is exactly where the obvious route fails:
deleting the non-Euler part of the arithmetic side leaves a form whose least eigenvalue is
negative at $17$ of $18$ configurations, including windows where the true form is
positive~\cite[\S15]{dhrep}.

\section{Origin of this result}\label{sec:origin}

\tier{T3} \emph{The trigger.} The forward-and-reverse question was raised immediately after
the companion note~\cite{dhnote} was deposited: a certified negative witness for one
function invites the question of whether a source-defined deformation could run the
implication the other way for another.

\emph{The author's prior inputs.} The certified witness and the strip bound of~\cite{dhnote},
and the finite-window transfer machinery of~\cite{deltaL3}.

\emph{The literature inputs.} De Bruijn's strip theorem in the form and normalisation
of~\cite[Thm.~7]{NW}; the existence of $\Lambda$ and $\Lambda\ge0$ from~\cite{RT};
the heat-flow normalisation and the improved upper bound from~\cite{P15}; the
Davenport--Heilbronn function from~\cite{DH36,BG11}.

\emph{The AI role.} AI assistance was used for the kernel derivation and its checking, the
ball-arithmetic certification, the adversarial controls and the drafting. An external review
raised six points, all accepted; four changed the mathematics rather than the wording, and
all six are logged in Appendix~\ref{app:review}. Conceptual direction, interpretation and
final responsibility are the author's.

\section{What this rules out, and what it does not}\label{sec:scope}

\tier{T3} \emph{Ruled out.} Any backward-positivity argument whose hypotheses are a subset
of $\{$conjugate reflection, rapid decay, bounded initial strip, kernel positivity, heat
evolution$\}$: such an argument, applied to Davenport--Heilbronn, would contradict the
certified witness. Also ruled out, in the weaker sense of being measured and found wanting:
the hope that locating a positive structure somewhere in the source --- for instance the
Bochner positivity \eqref{eq:bochner} --- is by itself progress toward positivity of the
zero form.

\emph{Not ruled out.} That intrinsic arithmetic dynamics can help. This note constrains one
deformation controlled by one package of general properties; it says nothing against a
deformation whose control uses the arithmetic. That $\LamDH$ is small, or large, within
\eqref{eq:bracket} --- the bracket is four orders of magnitude wide on a logarithmic scale
at the bottom and says nothing about where the transition is. That a prime-side
representation exists for $\tau\ne0$; none has been established, which is not the same as
none existing. And nothing whatever about $\Lambda$ for $\zeta$, or about RH: the forward
statements hold on every finite window and are silent at $\tau=0$, which is where the
arithmetic lives.

\emph{The methodological residue}, offered as judgement and not as a result: reduce the
admissible class of arithmetic objects first, using genuine arithmetic constraints, and
apply a general deformation theorem only to what survives. This note did the opposite ---
applied a general theorem to an unreduced class --- and the Davenport--Heilbronn control is
precisely what an unreduced class admits.

\appendix

\section{The certificate}\label{app:cert}

\tier{T2-C} All quantities from \src{data/lambda\_bracket.csv} and
\src{data/lambda\_bracket.txt}, produced by \src{src/lambda\_lower\_bound.py} in ball
arithmetic~\cite{Arb} at $400$ bits. One-sided bounds are plain decimals rounded to the safe
side; the balls they came from are printed alongside in the transcript.

\begin{center}
\begin{tabular}{@{}llr@{}}
\toprule
quantity & meaning & value \\
\midrule
$\operatorname{Re}z_1$ & disc centre, real part & $171.398696970755184343858535418$ \\
$\operatorname{Im}z_1$ & disc centre, imaginary part & $-0.617034364913274771106703921214$ \\
$r$ & disc radius & $0.3$ \\
$|\operatorname{Im}z_1|-r$ & margin from the real axis & $0.31703436491327477111$ \\
arcs & partition of $\partial D$ & $3000$ \\
$2r\sin(\Delta\theta/4)$ & arc ball radius & $0.000314159251004$ \\
--- & worst relative enclosure radius & $0.064946579$ \\
$m$ & $\min_{\partial D}|\HDH_0|$, floor & $6.5935981835929959399\times10^{-30}$ \\
$\frac{1}{2\pi}\sum\Delta\arg$ & winding number & $1$ \\
--- & worst increment $+$ uncertainties & $0.1313552700$ \\
$M$ & $\max_{\partial D}|\operatorname{Im}z|$ & $0.91703436491327477111$ \\
$\tau_0$ & perturbation radius & $0.1$ \\
$x_{\mathrm{cut}}$ & split point for $C$ & $2.000$ \\
$\int_0^{x_{\mathrm{cut}}}$ & head, rigorous integration & $0.048291904880526492445$ \\
$f(x_{\mathrm{cut}})$ & tail bound & $4.481699445\times10^{-809}$ \\
$C$ & ceiling & $0.048291904880526492446$ \\
$\varepsilon$ & $\min(\tau_0,m/(2C))$, floor & $6.8268151773112564447\times10^{-29}$ \\
\bottomrule
\end{tabular}
\end{center}

\section{Review response, and the retraction log}\label{app:retr}\label{app:review}

\tier{T3} An external review checked the cited theorem's hypotheses and parameter
convention and confirmed both, while declining to rerun the scripts or re-audit the kernel
construction step by step. It raised six points; all six were accepted, and four changed the
mathematics or its presentation rather than only the wording.

\begin{enumerate}\itemsep3pt
\item \textbf{The uniform positivity bound was presented wrongly.} The unfactored difference
tends to $0$, so the number quoted for it was not a uniform lower bound. The factored
form~\eqref{eq:pos} has an increasing bracket, $0.6763199240515301$ at $T=\pi/5$. \S2 of the
record and the corresponding block were rewritten; the review's correction was right and the
presentation was wrong.

\item \textbf{``No such inequality exists'' overstates it.} The correct statement is that no
\emph{universally applicable} backward argument on the shared hypotheses can work, and that
somewhere on the path the inequality, the coefficient control, the differentiability or
$B_\tau\succ0$ must fail. Adopted, and used in \S\ref{sec:transfer}.

\item \textbf{The slogan ``the deformation is informative exactly where the arithmetic is
absent'' is too strong}, since $\Phi_{F,\tau}$ still contains the $a_n$. Withdrawn and
replaced; see \S\ref{sec:transfer}.

\item \textbf{``There is no prime side at all'' should be ``no appropriate prime-side
representation has been established''.} Changed throughout.

\item \textbf{The toy polynomial is not a Fourier transform of an integrable rapidly
decaying kernel}, because such a transform is bounded on the real axis and a quadratic is
not. It is a polynomial heat-flow control reproducing the strip-shrinking rate, not an
example satisfying every hypothesis of Theorem~\ref{thm:dB}. Rewritten with the
boundedness argument, which is a better reason than the one originally given.

\item \textbf{A two-sided bracket is available}, by persistence of nonreal zeros for small
$|\tau|$, without differentiating the complete family; and positive kernels force a positive
semidefinite Fourier-correlation matrix, which coexists with a negative zero-based witness.
Both adopted: the first became Theorem~\ref{thm:lower} after being made quantitative in ball
arithmetic, the second is \S\ref{sec:filter}. Both observations are the review's and are
recorded as such.
\end{enumerate}

\noindent\textbf{Further corrections made while preparing this note.}

\begin{enumerate}\itemsep3pt
\setcounter{enumi}{6}
\item \textbf{A one-sided bound printed with round-to-nearest.} The record first gave
$m\ge6.5935981835929959400\times10^{-30}$, which is the enclosure's midpoint rounded to
nearest and therefore asserts about $4.5\times10^{-50}$ more than the enclosure supports.
Numerically harmless, wrong in principle for a certified floor. One-sided bounds are now
printed by rounding the ball's own endpoint to the safe side, giving the $m$ and $C$ values
of Appendix~\ref{app:cert}. \textbf{$\varepsilon$ is unaffected}, because it was already
computed from midpoint minus radius. Found by asking whether each published digit of a
one-sided bound was actually implied by the enclosure it came from.

\item \textbf{``rank 2'' where the numerics said rank 1}, in the positive semidefiniteness
lemma of \S\ref{sec:transfer}. In a pure parity sector one of the two rank-one terms
vanishes identically. The numerics caught the prose; the conclusion is unaffected.

\item \textbf{A cancellation artefact recorded as a residual.} An evenness check first
reported a large ``worst relative residual'' that was a low-precision evaluation of a
quantity needing roughly $105$ digits. The check now sets its precision per point and
reports an undecided point as undecided rather than as a failure; all points resolve, the
worst at $2.25262\times10^{-60}$.

\item \textbf{A patch that did not land.} A correction was first applied by a background
command that completed \emph{after} the run it was meant to fix, so the regenerated
transcript still carried the superseded text. Caught by searching the new transcript for the
stale string rather than assuming the edit had taken effect. This happened twice.
\end{enumerate}

\section{Reproduction}\label{app:repro}

\noindent \texttt{make verify}, in the note directory, re-derives every number printed above
from the audit's data files and fails on any disagreement, any undeclared literal, any
overstatement of precision relative to the source's own significant digits, or an unmet
prerequisite: it refuses to pass unless the certificate of Appendix~\ref{app:cert} is
present and its regression fixture passes. \texttt{make} builds the PDF.

\begin{center}
\begin{tabular}{@{}ll@{}}
\toprule
file & produced by \\
\midrule
\src{data/lambda\_bracket.csv}, \src{.txt} & \src{src/lambda\_lower\_bound.py --arcs 3000 --prec 400} \\
\src{data/validation.csv}, \src{audit\_log.txt} & \src{src/source\_flow\_audit.py --dps 40} \\
--- & \src{tests/test\_identities.py} \\
--- & \src{tests/test\_lambda\_certificate.py} \\
\src{note/lambda\_note.pdf} & \src{make -C note} \\
\bottomrule
\end{tabular}
\end{center}

\noindent The full record, with the feasibility gate, the adversarial controls and the
blocks not drawn on here, is the audit's \src{REPORT.md}~\cite{audit}; the sources actually
read, and which statement came from where, are in its \src{SOURCE\_AUDIT.md}.

\section*{Licence}
\addcontentsline{toc}{section}{Licence}

As the parent repository: code MIT, this note and the data CC~BY~4.0. CC~BY~4.0 requires
attribution; if you use the note or the data, please cite it by its DOI and the parent
paper~\cite{deltaL3} separately.

\section*{Acknowledgments}
\addcontentsline{toc}{section}{Acknowledgments}

AI assistance was used for algebra checking, numerical verification, adversarial review,
literature organization, and manuscript drafting. All conceptual development, research
direction, interpretation, and final responsibility for the content belong to the author.

\begin{thebibliography}{99}

\bibitem{BG11} E. Bombieri and A. Ghosh,
\emph{Around the Davenport--Heilbronn function},
Russian Math. Surveys \textbf{66} (2011), no.~2, 221--270;
Uspekhi Mat. Nauk \textbf{66} (2011), no.~2, 15--66.

\bibitem{CSV94} G. Csordas, W. Smith and R.~S. Varga,
\emph{Lehmer pairs of zeros, the de Bruijn--Newman constant $\Lambda$, and the Riemann
Hypothesis}, Constr. Approx. \textbf{10} (1994), no.~1, 107--129.
\textbf{Not consulted.} Cited only for the route the $\zeta$ lower-bound literature took,
as that route is described in~\cite[\S2.6]{NW}; nothing here uses a result of it.

\bibitem{DH36} H. Davenport and H. Heilbronn,
\emph{On the zeros of certain Dirichlet series},
J. London Math. Soc. \textbf{11} (1936), 181--185; \textbf{11} (1936), 307--312.

\bibitem{DB50} N.~G. de Bruijn,
\emph{The roots of trigonometric integrals},
Duke Math. J. \textbf{17} (1950), no.~3, 197--226; Theorem~13.
\textbf{Not consulted directly.} The theorem number is the one Newman gives
(\cite[p.~246]{Ne76}: ``due to de Bruijn [B, Theorem~13]''); the statement used here is
Newman's Theorem~1, which was read. The review MR0037351 confirms the strip result in
outline and records de Bruijn's own normalisation $\Xi(2z)=\int\varphi(t)e^{izt}dt$.

\bibitem{dhnote} E. De Jes\'us,
\emph{A Certified Negative Finite-Window Weil Witness for the Davenport--Heilbronn
Function, with the Remaining-Zero Form Positive}, Zenodo, 2026.
\href{https://doi.org/10.5281/zenodo.23608288}{doi:10.5281/zenodo.23608288}.

\bibitem{siegelnote} E. De Jes\'us,
\emph{Parity Blindness of Finite-Window Weil Detectors at a Real Zero}, Zenodo, 2026.
\href{https://doi.org/10.5281/zenodo.23623473}{doi:10.5281/zenodo.23623473}.

\bibitem{deltaL3} E. De Jes\'us,
\emph{A $\delta^2$--$L^3$ Sensitivity Law for Finite-Window Weil Positivity, with a
Function-Field Control of the Connes--van Suijlekom Pipeline}, Zenodo, 2026.
\href{https://doi.org/10.5281/zenodo.23250932}{doi:10.5281/zenodo.23250932}.

\bibitem{audit} E. De Jes\'us,
\emph{Source-defined deformation $\to$ fixed-window detector: a feasibility audit},
in \texttt{weil-window-control}, \texttt{extensions/source-flow-detector-audit-2026-10-24},
2026.

\bibitem{dhrep} E. De Jes\'us,
\emph{Davenport--Heilbronn control for the finite-window Weil detector},
in \texttt{weil-window-control}, \texttt{extensions/davenport-heilbronn-2026-10-23}, 2026.

\bibitem{Arb} F. Johansson,
\emph{Arb: Efficient arbitrary-precision midpoint-radius interval arithmetic},
IEEE Trans. Comput. \textbf{66} (2017), no.~8, 1281--1292.

\bibitem{Ne76} C.~M. Newman,
\emph{Fourier transforms with only real zeros},
Proc. Amer. Math. Soc. \textbf{61} (1976), no.~2, 245--251.
\textbf{Read.} Theorem~1 is the form of Theorem~\ref{thm:dB} used here, in the multiplier
convention used here; p.~246 carries both the attribution to de Bruijn and the
$\tfrac{1}{8}$ value in Newman's own $z$-variable.

\bibitem{NW} C.~M. Newman and W. Wu,
\emph{Constants of de Bruijn--Newman type in analytic number theory and statistical
physics}, arXiv:1901.06596v2. \textbf{Read.} The source for the hypotheses of
Theorem~\ref{thm:dB} and for the framework's positivity assumption; its Theorem~7 is the
same theorem in the multiplier convention $e^{\lambda t^2/2}$.

\bibitem{P15} D.~H.~J. Polymath,
\emph{Effective approximation of heat flow evolution of the Riemann $\xi$ function, and a
new upper bound for the de Bruijn--Newman constant}, arXiv:1904.12438v2.

\bibitem{RT} B. Rodgers and T. Tao,
\emph{The de Bruijn--Newman constant is non-negative},
Forum Math. Pi \textbf{8} (2020), e6, doi:10.1017/fmp.2020.6; arXiv:1801.05914v5.
Abstract and definitions read; used only for $\Lambda\ge0$ as a remark.

\end{thebibliography}

\end{document}
